% !TEX program = xelatex
% !TeX encoding = UTF-8
\documentclass[12pt,a4paper]{article}
\usepackage{fontspec}
\usepackage{polyglossia}
\setmainlanguage{russian}
\setotherlanguage{english}

% Шрифты для русского и латиницы
\setmainfont{Times New Roman}[Ligatures=TeX]
\setsansfont{Arial}
\setmonofont{Consolas}[Scale=MatchLowercase]

% Пакеты математики и форматирования
\usepackage{amsmath,amssymb,amsthm,mathtools}
\usepackage{geometry}
\geometry{left=2.5cm,right=2.5cm,top=2.5cm,bottom=2.5cm}
\usepackage{booktabs}
\usepackage{hyperref}
\usepackage{url}
\usepackage{enumitem}
\usepackage{float}
\usepackage{xcolor}
\usepackage{titlesec}
\usepackage{fancyhdr}

% Определение цветов
\definecolor{gold}{RGB}{255,215,0}
\definecolor{darkblue}{RGB}{0,0,139}
\definecolor{skyblue}{RGB}{0,102,204}

% Макросы YUCT
\newcommand{\Keff}{K_{\mathrm{eff}}}
\newcommand{\kappac}{\kappa_c}
\newcommand{\eps}{\varepsilon}
\newcommand{\df}{d_{\!f}}
\newcommand{\dint}{\ensuremath{D\!+\!I\!\cdot\!R}}
\newcommand{\Sodd}{S_{\mathrm{odd}}}
\newcommand{\Seven}{S_{\mathrm{even}}}

% Русские теоремы
\newtheorem{theorem}{Теорема}[section]
\newtheorem{definition}[theorem]{Определение}
\newtheorem{proposition}[theorem]{Предложение}
\newtheorem{conjecture}[theorem]{Гипотеза}
\theoremstyle{remark}
\newtheorem{remark}[theorem]{Замечание}

\hypersetup{
	colorlinks=true,
	linkcolor=blue,
	citecolor=blue,
	urlcolor=blue,
}

% Форматирование заголовков
\titleformat{\section}{\LARGE\bfseries\color{darkblue}}{\arabic{section}}{1em}{}
\titleformat{\subsection}{\Large\bfseries\color{darkblue}}{\arabic{subsection}}{1em}{}

% Нижний колонтитул
\fancypagestyle{firstpage}{%
	\fancyhf{}%
	\renewcommand{\headrulewidth}{0pt}%
	\renewcommand{\footrulewidth}{0pt}%
	\fancyfoot[C]{%
		\begin{minipage}[t]{\textwidth}
			\centering
			\textcolor{gold}{\rule{\textwidth}{1.6pt}}\\[5pt]
			{\small \textsuperscript{1}Исследования Якушева, \url{https://yuct.org/}} \hfill
			{\small Протокол YPSDC: \url{https://ypsdc.com/}}\\[-2pt]
			\textcolor{gold}{\rule{\textwidth}{1.6pt}}\\[5pt]
			\textbf{Единая теория координации Якушева 2026} \hfill \thepage\\[10pt]
		\end{minipage}%
	}%
}

\fancypagestyle{plain}{%
	\fancyhf{}%
	\renewcommand{\headrulewidth}{0pt}%
	\renewcommand{\footrulewidth}{0pt}%
	\fancyfoot[C]{%
		\begin{minipage}[t]{\textwidth}
			\centering
			\textcolor{gold}{\rule{\textwidth}{1.6pt}}\\[4pt]
			\textbf{Единая теория координации Якушева 2026} \hfill \thepage\\[4pt]
		\end{minipage}%
	}%
}

\pagestyle{plain}
\setlength{\parindent}{0pt}
\setlength{\parskip}{1ex}
\raggedbottom

\makeatletter
\let\@ensure@LTR\relax
\makeatother

% ============================================================
% Начало документа
% ============================================================
\begin{document}
	
	\begin{titlepage}
		\centering
		\vspace*{1cm}
		
		\Huge\textbf{Приложение F (пересмотренное): Фрактальная динамика экономических словарей и теория координации денег}
		
		\vspace{0.5cm}
		\Large\textbf{Теоретическая основа для обменных курсов, валютных зон и самоорганизующихся экономических границ}
		
		\vspace{1.5cm}
		
		\Large\textbf{Алексей В. Якушев\textsuperscript{1}}
		
		Теория YUCT: \url{https://yuct.org/}\\
		Протокол YPSDC: \url{https://ypsdc.com/}
		
		\vspace{2cm}
		\textbf{DOI:} \url{https://doi.org/10.5281/zenodo.18444598}
		
		\vspace{1cm}
		
		\large
		Июнь 2026 \\ \large Версия 3.0 (заменяет предыдущее Приложение F)
		
		\vspace{2cm}
		
		\begin{minipage}{0.9\textwidth}
			\begin{abstract}
				\noindent
				В данном приложении предлагается коренной пересмотр подхода YUCT к экономике. Мы отказываемся от всех предыдущих попыток эмпирической подгонки ВВП или волатильности рынка с универсальным показателем \(\beta = 0.67\) и вместо этого разрабатываем теоретическую основу из первых принципов, опирающуюся на онтологию координации. Экономика моделируется как сеть взаимодействующих \textbf{экономических словарей} (фирм, секторов, стран), каждый из которых характеризуется эффективностью координации \(K_{\mathrm{eff}}\). Мы показываем, что граница между любыми двумя словарями демонстрирует фрактальную динамику, аналогичную множеству Мандельброта: самоподобную, сложную и управляемую универсальным законом ошибки \(\varepsilon = \kappa_c \alpha K_{\mathrm{eff}}^{-\beta}\). Деньги определяются как \textbf{параметр координации}, регулирующий активацию через границы словарей; таким образом, обменные курсы определяются отношением эффективностей координации. Эта структура даёт несколько количественных, фальсифицируемых предсказаний: (1) Показатель Херста для высокочастотных рядов обменных курсов должен стремиться к \(H \approx 1/\beta \approx 1.5\) на больших временных масштабах; (2) Мультифрактальный спектр финансовых доходностей должен демонстрировать левостороннюю асимметрию универсальной формы, определяемой первичным алгебраическим циклом; (3) Волатильность \(\sigma\) обменного курса масштабируется с разрывом координации \(\Delta K_{\mathrm{eff}}\) между словарями как \(\sigma \propto (\Delta K_{\mathrm{eff}})^{-0.67}\). В данном приложении не производится никаких эмпирических подгонок; все предсказания сформулированы в проверяемой форме, как только станут доступны надёжные прокси для \(K_{\mathrm{eff}}\) в экономических системах. Эта работа превращает экономику YUCT из набора сомнительных корреляций в строгую, фальсифицируемую исследовательскую программу.
			\end{abstract}
		\end{minipage}
		
		\vspace{1cm}
		
		\noindent
		\small\textbf{Ключевые слова:} YUCT, эффективность координации, экономические словари, обменные курсы, фрактальные границы, множество Мандельброта, показатель Херста, мультифрактальный анализ, деньги как параметр координации.
		
		\vspace{0.5cm}
		
		\noindent
		\textcopyright\ 2026 Исследования Якушева. Все права защищены.
	\end{titlepage}
	
	\tableofcontents
	\newpage
	
	% ============================================================
	% РАЗДЕЛ 1: ВВЕДЕНИЕ И МОТИВАЦИЯ
	% ============================================================
	\section{Введение: почему необходим новый старт}
	
	В предыдущих версиях этого приложения предпринимались попытки доказать, что макроэкономические временные ряды — волатильность роста ВВП, ошибки прогнозирования криптовалют, индикаторы геополитического риска — следуют универсальному закону масштабирования \(\varepsilon \propto K_{\mathrm{eff}}^{-0.67}\) теории YUCT. Последующая независимая проверка показала, что многие из упомянутых таблиц и регрессий были либо ошибочными, либо основаны на постфактумном отборе благоприятных случаев. Данная редакция полностью признаёт эти недостатки и отзывает все предыдущие эмпирические утверждения.
	
	Основная ошибка заключалась не в самом фреймворке YUCT, а в попытке навязать прямую эмпирическую подгонку без предварительной разработки адекватной теоретической модели того, что же такое «экономический словарь» и как можно измерить его эффективность координации. Экономика — это не физика частиц: экономические агенты не являются точечными частицами, и их массы не являются фундаментальными константами. Поэтому поиск простой массовой шкалы в экономических данных был изначально ошибочным.
	
	В данном приложении предлагается принципиально иной подход. Вместо того чтобы рассматривать экономику как совокупность объектов, размеры которых должны попадать на дискретный спектр, мы трактуем её как \textbf{динамическую сеть взаимодействующих словарей}, границы которых являются основным местом возникновения интересных явлений. Обменные курсы, процентные ставки и финансовые волатильности — это не свойства изолированных агентов, а \textbf{эмерджентные свойства границ} между словарями. Геометрия этих границ, как мы утверждаем, фрактальна и самоподобна и управляется тем же универсальным показателем \(\beta = 0.67\), который проявляется во всех других областях YUCT.
	
	В этом приложении не производится никаких эмпирических подгонок. Все количественные утверждения выводятся математически из аксиом YUCT и представляются как фальсифицируемые предсказания для будущих исследований.
	
	% ============================================================
	% РАЗДЕЛ 2: ЭКОНОМИЧЕСКИЕ СЛОВАРИ И ИХ ГРАНИЦЫ
	% ============================================================
	\section{Экономические словари и их границы}
	
	\subsection{Определение экономического словаря}
	
	\begin{definition}[Экономический словарь]
		\textbf{Экономический словарь} \(\mathcal{D}_i\) — это структурированное множество возможных действий, контрактов, протоколов производства, правил ценообразования и институциональных норм, которые экономический агент (фирма, домохозяйство, сектор, государство) может активировать в ответ на внешние индикаторы. Размер словаря измеряется энтропией Шеннона \(H(\mathcal{D}_i)\).
	\end{definition}
	
	Примеры:
	\begin{itemize}
		\item Словарь фирмы включает её каталог продуктов, контракты на поставку, алгоритмы ценообразования и внутренние управленческие процедуры.
		\item Словарь национальной экономики включает её правовую систему, налоговое законодательство, торговые соглашения, рамки денежно-кредитной политики и финансовую инфраструктуру.
	\end{itemize}
	
	Эффективность координации словаря \(i\) равна
	\begin{equation}
		K_{\mathrm{eff}}^{(i)} = \frac{H(\mathcal{A}_i)}{H(\kappa_i)},
		\label{eq:keff_econ}
	\end{equation}
	где \(\mathcal{A}_i\) — множество действий, которые может породить словарь, а \(\kappa_i\) — индикатор (ценовой сигнал, регуляторное предписание, рыночный приказ), который их запускает.
	
	\subsection{Граница между двумя словарями}
	
	\begin{definition}[Граница словарей]
		\textbf{Граница} \(\partial\mathcal{D}_{ij}\) между двумя экономическими словарями \(\mathcal{D}_i\) и \(\mathcal{D}_j\) — это множество всех транзакций, контрактов и информационных обменов, в которых индикатор, возникший в одном из словарей, активирует вход в другом.
	\end{definition}
	
	Граница является естественным местом для:
	\begin{itemize}
		\item обменных курсов (для национальных словарей),
		\item транзакционных издержек и ценовых разрывов (для словарей на уровне фирм),
		\item информационной асимметрии и неблагоприятного отбора.
	\end{itemize}
	
	\begin{proposition}[Параметр координации на границе]
		Эффективность координации через границу \(\partial\mathcal{D}_{ij}\) измеряется безразмерным \textbf{параметром координации} \(\chi_{ij}\), определяемым как
		\begin{equation}
			\chi_{ij} = \frac{K_{\mathrm{eff}}^{(j)}}{K_{\mathrm{eff}}^{(i)}} \cdot \rho_{ij},
			\label{eq:chi}
		\end{equation}
		где \(\rho_{ij}\) — резонансный фактор, определяющий совместимость двух словарей (общие стандарты, общий язык, правовая гармонизация и т.д.).
	\end{proposition}
	
	\textbf{Деньги как параметр координации.} В случае двух национальных экономик с независимыми валютами параметр координации \(\chi_{ij}\) является в точности \textbf{обменным курсом} (единиц валюты \(j\) за единицу валюты \(i\)) с точностью до нормировочной константы. Более высокое \(K_{\mathrm{eff}}^{(j)}\) по отношению к \(K_{\mathrm{eff}}^{(i)}\) означает, что словарь \(j\) более «ценен» — его входы активируются более надёжно — и поэтому его валюта требует премии.
	
	% ============================================================
	% РАЗДЕЛ 3: ФРАКТАЛЬНАЯ ДИНАМИКА ГРАНИЦЫ
	% ============================================================
	\section{Фрактальная динамика границы}
	
	\subsection{Аналогия с множеством Мандельброта}
	
	Множество Мандельброта \(\mathcal{M}\) определяется итеративным отображением \(z_{n+1} = z_n^2 + c\) с \(z_0 = 0\). Точки \(c\), для которых последовательность остаётся ограниченной, принадлежат \(\mathcal{M}\); точки, где она расходится, лежат вне его. Граница \(\partial\mathcal{M}\) фрактальна, бесконечно сложна и демонстрирует самоподобие на всех масштабах.
	
	Мы предлагаем экономический аналог. Пусть два словаря \(\mathcal{D}_i\) и \(\mathcal{D}_j\) взаимодействуют через последовательность транзакций, индексированных дискретными временными шагами \(n\). Состояние координации после \(n\) транзакций описывается комплексной переменной \(z_n \in \mathbb{C}\), действительная и мнимая части которой представляют степень сходства в двух ортогональных измерениях (например, синхронизация цен и выполнение контрактов). Правило обновления имеет вид
	\begin{equation}
		z_{n+1} = z_n^2 + \chi_{ij},
		\label{eq:mandel_econ}
	\end{equation}
	где \(\chi_{ij} \in \mathbb{C}\) — (комплексный) параметр координации из уравнения~(\ref{eq:chi}). Модуль \(|\chi_{ij}|\) измеряет силу связи, а его фаза — фазовый сдвиг между циклами активации двух словарей.
	
	\begin{conjecture}[Экономическое соответствие Мандельброта]
		Пара словарей \((\mathcal{D}_i, \mathcal{D}_j)\) может поддерживать ограниченную и устойчивую координацию (т.е. последовательность \(\{z_n\}\) остаётся ограниченной) тогда и только тогда, когда \(\chi_{ij}\) лежит внутри обобщённого множества Мандельброта \(\mathcal{M}_{\mathrm{econ}}\). Граница \(\partial\mathcal{M}_{\mathrm{econ}}\) отделяет области устойчивой координации (стабильные обменные курсы, предсказуемые торговые потоки) от областей хаотической дивергенции (валютные кризисы, гиперинфляция, коллапс торговли). Геометрия этой границы фрактальна, с самоподобными структурами на всех временных масштабах.
	\end{conjecture}
	
	\subsection{Универсальное масштабирование на границе}
	
	Поскольку динамика на границе подчиняется тому же универсальному закону ошибки координации, который действует повсюду в YUCT, статистические свойства обменных курсов вблизи границы должны подчиняться универсальным законам масштабирования. Пусть \(\sigma_{ij}(\Delta t)\) — волатильность обменного курса между словарями \(i\) и \(j\), измеренная на временном горизонте \(\Delta t\). Тогда из универсального закона ошибки \(\varepsilon \propto K_{\mathrm{eff}}^{-\beta}\) получаем:
	
	\begin{proposition}[Масштабирование волатильности обменного курса]
		Волатильность параметра координации \(\chi_{ij}\) масштабируется с разрывом координации \(\Delta K_{\mathrm{eff}} = |K_{\mathrm{eff}}^{(i)} - K_{\mathrm{eff}}^{(j)}|\) как
		\begin{equation}
			\sigma(\chi_{ij}) \propto (\Delta K_{\mathrm{eff}})^{-\beta} = (\Delta K_{\mathrm{eff}})^{-0.67}.
			\label{eq:vol_scaling}
		\end{equation}
		Пары словарей с очень близкими эффективностями координации (малый разрыв \(\Delta K_{\mathrm{eff}}\)) демонстрируют высокую, хаотическую волатильность; пары с большим разрывом показывают низкую, стабильную волатильность.
	\end{proposition}
	
	Это фальсифицируемое предсказание. Как только будут разработаны надёжные прокси для \(K_{\mathrm{eff}}^{(i)}\) (например, из индикаторов качества институтов, сложности торговли или информационной энтропии законов), уравнение~(\ref{eq:vol_scaling}) может быть проверено на панельных данных по странам.
	
	\subsection{Фрактальные свойства финансовых временных рядов}
	
	Известным эмпирическим фактом в эконофизике является то, что финансовые временные ряды обладают мультифрактальными свойствами: кластеризация волатильности, долговременная память и распределения доходностей с тяжёлыми хвостами. В рамках YUCT это не случайности, а необходимые следствия фрактальной динамики на границах словарей.
	
	\begin{proposition}[Показатель Херста на границе]
		Для пары словарей, работающей вблизи критической границы устойчивой координации, показатель Херста \(H\) ряда обменного курса удовлетворяет
		\begin{equation}
			H \to \frac{1}{\beta} = \frac{3}{2} \quad \text{при} \quad \Delta t \to \infty.
			\label{eq:hurst}
		\end{equation}
		На более коротких временных масштабах \(H\) демонстрирует переход к меньшим значениям, что отражает мультифрактальную природу границы.
	\end{proposition}
	
	Это предсказание может быть проверено с помощью стандартного R/S-анализа или DFA на высокочастотных данных обменных курсов. Значение \(H = 1.5\) характерно для процесса с очень долгой памятью, близкого к сингулярности — именно такое поведение наблюдается перед крупными финансовыми кризисами (например, кризисом 2008 года).
	
	\subsection{Внутренняя фрактальная организация одного словаря}
	
	В то время как границы между словарями фрактальны и хаотичны, внутренняя структура хорошо скоординированного словаря демонстрирует дискретную иерархическую упорядоченность. Записи словаря распределены не равномерно, а образуют самоподобное ветвящееся дерево правил, контрактов и протоколов. Размеры фирм, распределение доходов и иерархия городских систем — всё это проявления этой внутренней фрактальной структуры.
	
	В отличие от предыдущей неудачной попытки поместить целые экономики на массовую шкалу, мы теперь предсказываем, что \textbf{внутри} одного словаря размеры подразделений (фирмы, города, классы доходов) должны следовать степенному закону, показатель которого связан с фрактальной размерностью \(d_f\) и универсальным \(\beta\). Известный закон Ципфа для городов (показатель \(\approx 1\)) и закон Парето для доходов (показатель \(\approx 1.5\)) являются не изолированными курьёзами, а различными проекциями одной и той же лежащей в основе координационной геометрии.
	
	% ============================================================
	% РАЗДЕЛ 4: МАТЕМАТИЧЕСКАЯ ОСНОВА
	% ============================================================
	\section{Математическая основа: лагранжиан координации на границе}
	
	Чтобы сделать изложенную выше качественную картину математически строгой, мы встраиваем её в 19-мерный лагранжев формализм YUCT. Пусть экономический сектор описывается полем \(\Psi_{72}(x^\mu, y^m)\), где \(x^\mu\) — четыре координаты пространства-времени, а \(y^m\) — 15 внутренних (компактных) координат. Словарь \(\mathcal{D}_i\) соответствует определённому вакуумному ожиданию \(\langle \Psi_{72} \rangle_i\), а граница между двумя словарями — это решение в виде доменной стенки, интерполирующей между \(\langle \Psi_{72} \rangle_i\) и \(\langle \Psi_{72} \rangle_j\).
	
	Эффективное действие на границе имеет вид
	\begin{equation}
		S_{\mathrm{boundary}} = \int d^4x \sqrt{-g} \left[ \frac{1}{2} \kappa (\nabla \chi)^2 + V(\chi) \right],
		\label{eq:action}
	\end{equation}
	где \(\chi(x)\) — локальный параметр координации (поле обменного курса), \(\kappa\) — жёсткость границы, а потенциал \(V(\chi)\) задаётся выражением
	\begin{equation}
		V(\chi) = V_0 \left( e^{\lambda(\chi - \chi_0)} - \lambda(\chi - \chi_0) - 1 \right), \qquad \lambda = \frac{1}{\beta} = \frac{3}{2}.
		\label{eq:potential}
	\end{equation}
	Это тот же потенциал, который управляет скалярным полем координации в гравитационном секторе (Приложение G), что обеспечивает универсальность.
	
	Классическое уравнение движения для \(\chi\) имеет вид
	\begin{equation}
		\kappa \Box \chi - V'(\chi) = 0,
		\label{eq:eom}
	\end{equation}
	а его статистические флуктуации удовлетворяют
	\begin{equation}
		\langle (\delta \chi)^2 \rangle \propto \int \frac{d^4k}{(2\pi)^4} \frac{1}{\kappa k^2 + V''(\chi_0)} \propto \chi_0^{-2\beta},
		\label{eq:fluctuations}
	\end{equation}
	что воспроизводит закон масштабирования волатильности из уравнения~(\ref{eq:vol_scaling}).
	
	% ============================================================
	% РАЗДЕЛ 5: ФАЛЬСИФИЦИРУЕМЫЕ ПРЕДСКАЗАНИЯ
	% ============================================================
	\section{Фальсифицируемые предсказания}
	
	Разработанная выше структура даёт несколько количественных предсказаний, которые могут быть проверены без какой-либо подгонки данных:
	
	\begin{enumerate}[label=(\arabic*)]
		\item \textbf{Показатель Херста обменных курсов.} Для основных свободно плавающих валютных пар (например, EUR/USD, USD/JPY) долгосрочный показатель Херста, рассчитанный по дневным данным за 20 лет, должен лежать в диапазоне \(1.3 \le H \le 1.7\). Нулевая гипотеза мейнстримной финансовой теории (случайное блуждание, \(H = 0.5\)) должна быть отвергнута с высокой достоверностью.
		
		\item \textbf{Масштабирование волатильность–разрыв координации.} Как только будет построен надёжный прокси для \(K_{\mathrm{eff}}\) (например, из индикаторов качества управления Всемирного банка в сочетании с данными о сложности торговли), поперечная волатильность обменных курсов должна масштабироваться с абсолютной разностью \(\Delta K_{\mathrm{eff}}\) как \(\sigma \propto (\Delta K_{\mathrm{eff}})^{-0.67}\). Показатель степени может быть оценён с помощью линейной регрессии в логарифмических координатах; предсказание состоит в том, что наклон составит \(-0.67 \pm 0.10\).
		
		\item \textbf{Асимметрия мультифрактального спектра.} Спектр сингулярности \(f(\alpha)\) доходностей обменных курсов должен демонстрировать левостороннюю асимметрию с коэффициентом асимметрии \(A_\alpha \approx 0.30 \pm 0.05\), что отражает доминирование крупных скоординированных обвалов над крупными скоординированными подъёмами. Это может быть проверено с помощью стандартных программ MF-DFA на общедоступных данных.
		
		\item \textbf{Критическое замедление перед валютными кризисами.} Когда валютная пара приближается к критической границе (т.е. когда \(\Delta K_{\mathrm{eff}} \to 0\)), время автокорреляции волатильности должно расходиться. Это предсказывает, что перед крупным валютным кризисом должно наблюдаться устойчивое увеличение показателя Херста и сужение мультифрактального спектра, что даёт сигнал раннего предупреждения.
	\end{enumerate}
	
	\textbf{Статус предсказаний:} Все предсказания выводятся математически из аксиом YUCT без подгонки под экономические данные. Они представлены в форме, которая может быть непосредственно проверена независимыми исследователями. Подтверждение любого из них предоставило бы сильные доказательства в пользу теории координации денег; фальсификация всех четырёх потребовала бы существенного пересмотра框架.
	
	% ============================================================
	% РАЗДЕЛ 6: ОТКРЫТЫЕ ВОПРОСЫ И ПУТЬ ВПЕРЁД
	% ============================================================
	\section{Открытые вопросы и путь вперёд}
	
	\begin{enumerate}[label=(\roman*)]
		\item \textbf{Построение надёжного прокси для \(K_{\mathrm{eff}}\) стран.} Наиболее насущной задачей является разработка метода оценки эффективности координации национальной экономики на основе общедоступных данных. Кандидаты: энтропия законодательства, сложность экспортной корзины (индекс Идальго–Хаусмана), качество институтов (индикаторы управления Всемирного банка) и информационные меры микроструктуры рынка. Составной индекс, откалиброванный на небольшой группе референтных стран, затем может быть экстраполирован на всю панель.
		
		\item \textbf{Численное моделирование экономического отображения Мандельброта.} Отображение \(z_{n+1} = z_n^2 + \chi\) должно быть исследовано систематически. Для каких значений \(\chi\) последовательность остаётся ограниченной? Какова фрактальная размерность границы? Как влияет добавление шума (конечная температура) на структуру? Эти вопросы могут быть решены с помощью простых скриптов на Python и могут выявить неожиданные количественные связи с реальной динамикой обменных курсов.
		
		\item \textbf{Связь с устоявшейся эконофизикой.} Предсказания, перечисленные в разделе~5, значительно пересекаются с существующими результатами эконофизики (мультифрактальность финансовых доходностей, долгая память в волатильности, степенные хвосты). Систематический обзор этой литературы, интерпретированной через призму YUCT, позволил бы немедленно проверить框架 без сбора новых данных.
	\end{enumerate}
	
	% ============================================================
	% РАЗДЕЛ 7: ЗАКЛЮЧЕНИЕ
	% ============================================================
	\section{Заключение}
	
	Мы полностью реструктурировали подход YUCT к экономике. Новая структура отказывается от неудачной стратегии принудительного размещения экономических агрегатов на массовую шкалу и вместо этого развивает теорию из первых принципов об экономических словарях и их границах. Деньги определяются как параметр координации, регулирующий активацию через словари, а обменные курсы оказываются управляемыми тем же универсальным показателем \(\beta = 0.67\), который проявляется повсюду в YUCT. Полученные предсказания количественны, фальсифицируемы и независимы от какой-либо последующей подгонки данных. Это превращает экономику YUCT из обузы в многообещающую, строгую исследовательскую программу.
	
	\section*{Отказ от ответственности}
	
	Данное приложение является теоретической работой. Не делается никаких эмпирических утверждений, и оно не предназначено для использования в качестве инвестиционных или политических рекомендаций. Все количественные предсказания представлены для научной проверки и могут быть фальсифицированы будущими данными.
	
	\section*{Доступность данных}
	
	Для этого теоретического исследования не было создано новых данных. Все упомянутые источники являются общедоступными.
	
	\begin{thebibliography}{99}
		\bibitem{AppendixA} А.В. Якушев, ``Приложение A: 19-мерный лагранжиан YUCT,'' в \emph{YUCT}, Zenodo (2026).
		\bibitem{AppendixL} А.В. Якушев, ``Приложение L: Фрактальное измерение ошибки координации,'' в \emph{YUCT}, Zenodo (2026).
		\bibitem{AppendixG} А.В. Якушев, ``Приложение G: Гравитация как геометрическое натяжение,'' в \emph{YUCT}, Zenodo (2026).
		\bibitem{AppendixX} А.В. Якушев, ``Приложение X: Обобщение теории информации Шеннона,'' в \emph{YUCT}, Zenodo (2026).
		\bibitem{Mandelbrot1982} Б.Б. Мандельброт, \emph{Фрактальная геометрия природы}, W.H. Freeman (1982).
		\bibitem{Pareto1896} В. Парето, \emph{Курс политической экономии}, Rouge (1896).
		\bibitem{Zipf1949} Дж.К. Ципф, \emph{Поведение человека и принцип наименьшего усилия}, Addison-Wesley (1949).
		\bibitem{Hidalgo2009} С.А. Идальго, Р. Хаусман, ``Строительные блоки экономической сложности,'' \emph{PNAS} \textbf{106}, 10570–10575 (2009).
		\bibitem{Kantelhardt2002} Й.В. Кантельхардт и др., ``Мультифрактальный анализ флуктуаций с детрендом для нестационарных временных рядов,'' \emph{Physica A} \textbf{316}, 87–114 (2002).
	\end{thebibliography}
	
\end{document}